\begin{theorem}[5.5.8 (fórmula de las dimensiones)]
\label{IV.2.s5.5.8-fr}
Sean $A$ un anillo local noetheriano íntegro, $B$ un anillo íntegro que
contiene a $A$ y que es una $A$-álgebra de tipo finito, $\mathfrak q$ un ideal
primo de $B$ tal que $\mathfrak q\cap A$ es el ideal maximal $\mathfrak m$ de
$A$, y $k$ y $k'$ los cuerpos residuales de $A$ y $B_{\mathfrak q}$,
respectivamente. Se tiene entonces la desigualdad
\begin{equation}
\dim(A)+\operatorname{deg.tr}_A B
\geq\dim(B_{\mathfrak q})+\operatorname{deg.tr}_k k'.
\tag{5.5.8.1}\label{IV.2.s5.5.8.1-fr}
\end{equation}
Además, los dos miembros son iguales si, para toda sub-$A$-álgebra $A'$ de
tipo finito de $B[T]$ y todo ideal maximal $\mathfrak m'$ de $A'$ tal que
$\mathfrak m'\cap A=\mathfrak m$, $A'_{\mathfrak m'}$ es catenario.
\end{theorem}

\begin{proof}
Sea $B=A[x_1,\ldots,x_n]$ y razonemos por inducción sobre $n$. Pongamos
$C=A[x_1,\ldots,x_{n-1}]$ y $\mathfrak r=\mathfrak q\cap C$;
$C_{\mathfrak r}$ es un anillo local íntegro noetheriano, y si se pone
$S=C\setmin\mathfrak r$, $B'=S^{-1}B$, $\mathfrak q'=S^{-1}\mathfrak q$, se tiene
$B_{\mathfrak q}=B'_{\mathfrak q'}$, $B'=C_{\mathfrak r}[x_n]$ y
$\mathfrak q'\cap C_{\mathfrak r}=\mathfrak rC_{\mathfrak r}$; además, los
cuerpos de fracciones
\oldpage[IV-2]{97}
de $B'$ y $C_{\mathfrak r}$ son los de $B$ y $C$,
respectivamente. Si $k_1$ es el cuerpo residual de $C_{\mathfrak r}$, se
tiene, por tanto, según (5.5.7.1)

\begin{equation}
\dim(C_{\mathfrak r})+\operatorname{deg.tr}_C B
\geq\dim(B_{\mathfrak q})+\operatorname{deg.tr}_{k_1}k'.
\tag{5.5.8.2}\label{IV.2.s5.5.8.2-fr}
\end{equation}
Por otra parte, la hipótesis de inducción da
\begin{equation}
\dim(A)+\operatorname{deg.tr}_A C
\geq\dim(C_{\mathfrak r})+\operatorname{deg.tr}_k k_1
\tag{5.5.8.3}\label{IV.2.s5.5.8.3-fr}
\end{equation}
de donde se obtiene (5.5.8.1) sumando (5.5.8.2) y (5.5.8.3) miembro a
miembro. Para demostrar la segunda aserción, observemos primero que toda
sub-$A$-álgebra de tipo finito de $C[T]$ es también una sub-$A$-álgebra de
tipo finito de $B[T]$; por inducción sobre $n$, se puede, por tanto, suponer
que los dos miembros de (5.5.8.3) son iguales. Por otra parte, para ver que
los dos miembros de (5.5.8.2) son iguales, basta, en virtud de (5.5.7),
comprobar que, si $\mathfrak n$ es un ideal maximal de $C_{\mathfrak r}[T]$
tal que $\mathfrak n\cap C_{\mathfrak r}=\mathfrak rC_{\mathfrak r}$,
entonces $(C_{\mathfrak r}[T])_{\mathfrak n}$ es catenario; pero se tiene
$C_{\mathfrak r}[T]=S'^{-1}C[T]$, donde $S'=C\setmin\mathfrak r$; el ideal
$\mathfrak n$ es, por tanto, de la forma $S'^{-1}\mathfrak n'$, donde
$\mathfrak n'$ es un ideal primo de $C[T]$ tal que
$\mathfrak n'\cap C=\mathfrak r$, de donde
$(C_{\mathfrak r}[T])_{\mathfrak n}=(C[T])_{\mathfrak n'}$; si
$\mathfrak m'$ es un ideal maximal de $C[T]$ que contiene a $\mathfrak n'$,
$(C[T])_{\mathfrak n'}$ es, por tanto, un anillo local del anillo
$(C[T])_{\mathfrak m'}$, y como por hipótesis este último es catenario, lo
mismo ocurre con $(C[T])_{\mathfrak n'}$ (\textbf 0, 16.1.4).
\end{proof}

